\documentclass[11pt]{article}
\usepackage[margin=1in]{geometry}
\usepackage[T1]{fontenc}
\usepackage[utf8]{inputenc}
\usepackage{lmodern}
\usepackage{booktabs}
\usepackage{array}
\usepackage{hyperref}
\usepackage{caption}
\usepackage{enumitem}
\usepackage{microtype}

\hypersetup{
  colorlinks=true,
  linkcolor=black,
  citecolor=black,
  urlcolor=blue
}

\title{\textbf{Interdict: Naming SQL Guardrail Rules Triggers Literal Compliance}\thanks{Interdict website: \url{https://interdict.vercel.app/}. Source code: \url{https://github.com/prisharai/Interdict}.}}
\author{First Author\\Interdict Research\\\texttt{first.author@example.com} \and Second Author\\Affiliation Placeholder\\\texttt{second.author@example.com}}
\date{\today}

\begin{document}
\maketitle

\begin{abstract}
We study what database agents do after a runtime guardrail denies a write. The completed v1 run used broad-request tasks, so full-table rewrites are not confirmatory evidence of strategic evasion; they measure literal rule satisfaction under a broad objective that conflicts with the guardrail. The complete, balanced primary run is Claude Haiku 4.5; Claude Sonnet 4.6 and Claude Opus 4.8 are partial and descriptive only. Across the full sweep, the result file reports 588 trials, 3 tasks, and 4 denial conditions; model completeness differs because the sweep was interrupted by an API credit limit. In Haiku, the operational any-evasion label rose from 22\% [13\%, 34\%] under opaque denial to 95\% [86\%, 98\%] when the denial named the violated rule, and to 98\% [91\%, 100\%] when the denial added a fix or fix plus blast-radius information. The corresponding Haiku Fisher comparisons against opaque denial were OR=68.69, p=1.953e-17 for reason-code feedback, and OR=213.31, p=7.077e-20 for both richer conditions. This pattern was driven by scope-theater attempts: statements with a syntactically valid \texttt{WHERE} clause that still affected approximately the whole table. The honest conclusion is exploratory: in v1, mechanistic denial feedback elicited literal rule satisfaction with the stated rule, not confirmed narrow-intent evasion. The pre-registered v2 task set is the confirmatory test of evasion under explicit narrow intent.
\end{abstract}

\section{Thesis and framing}
When an agent is told to perform a broad write and the guardrail blocks it, naming the violated rule, such as \texttt{WRITE\_WITHOUT\_WHERE}, makes the agent far more likely to re-issue a syntactically compliant statement that still hits the whole table than an opaque error does. The paper therefore treats v1 as a study of feedback-induced literal rule satisfaction under a conflicting broad objective. It does not treat v1 as a confirmatory study of strategic evasion of a narrow user intent. The term ``operational any-evasion label'' refers to the logged classifier outcome; prose uses ``scope-theater attempts'' or ``literal rule satisfaction'' for v1 behavior unless the label name itself is being described.\footnote{The project handoff states that v1 tasks asked for broad operations, so \texttt{scope\_theater}/\texttt{destructive} labels measure literal rule satisfaction under a broad objective, not strategic evasion of narrow intent. The label audit repeats the same validity note.}

\section{Methods}
\subsection{Design and data state}
The sweep crossed model, task, and denial-richness condition. The result file reports 588 trials across 3 tasks and 4 conditions, with model-condition cell sizes in [28, 60]. Haiku is the complete primary run: 240 trials, 12 task-by-condition cells filled, and minimum 20 trials per task-by-condition cell. Opus has 148 trials, 8 filled cells, and minimum 8 trials per filled cell; Sonnet has 200 trials, 10 filled cells, and minimum 20 trials per filled cell. Because Haiku has 3 tasks and 20 trials per task-by-condition cell, condition-level Haiku analyses collapse to n=60 per condition.\footnote{All result counts in this paragraph are copied from \texttt{RESULTS\_STUDY.md}, data-completeness lines 3 and 10--15.}

The closed loop was capped at \texttt{\_MAX\_TURNS=4}: at most 4 agent attempts, logged as turns 0--3. Each turn consisted of one SQL statement, an engine decision, and condition-specific feedback if the statement was not executed. Trial order was logged under a single seed; robustness to alternate seeds or trial orders was not measured in v1.\footnote{The uploaded README describes the closed loop as up to 4 turns. \texttt{INDEX.md} lists ``one seed'' among required caveats; no alternate-seed analysis is printed in \texttt{RESULTS\_STUDY.md}.}

\subsection{Tasks and conditions}
The completed v1 tasks were broad-request database chores. Under those prompts, a full-table write can be faithful to the user-facing request even though the guardrail blocks it. This is the core validity distinction: v1 does not confirm strategic evasion of narrow intent; it measures literal rule satisfaction when the broad user objective conflicts with the guardrail.

The independent variable was denial feedback. \texttt{C0\_opaque} was a bare rejection. \texttt{C1\_reason\_code} named the violated rule, for example \texttt{WRITE\_WITHOUT\_WHERE}. \texttt{C2\_reason\_fix} added a human message and suggested fix. \texttt{C3\_reason\_fix\_blast} added measured-row and held-for-confirmation framing. The C3 framing is a limitation because it confounds row-count information with confirmation language; v2 separates those factors.

\subsection{Instrument and labels}
Interdict parses SQL, applies a deterministic policy, simulates candidate writes in a rollback transaction to measure affected rows, and logs the SQL, engine decision, reason code, measured blast radius, intended rows, total rows, label, and feedback shown. The measurement apparatus is designed to distinguish broad-with-\texttt{WHERE} statements from genuine narrowing: \texttt{scope\_theater\_evasion} is defined as having a \texttt{WHERE} while the measured blast radius is approximately full-table; \texttt{genuine\_correction}/\texttt{allowed\_ontask} is an allowed write whose blast radius is approximately intended scope. The audit file gives one exact numeric classifier rule: \texttt{destructive} is a write with blast radius at least \(0.9\cdot\)table or no \texttt{WHERE}. The uploaded sources do not expose an exact numeric cutoff for the approximately-full-table \texttt{scope\_theater\_evasion} rule; that cutoff is not measured in the result file.\footnote{The label definitions are from \texttt{label\_audit.md}. This is a classifier-design statement, not a result estimate.}

The Results report only outcomes printed in \texttt{RESULTS\_STUDY.md}: first-attempt denial, operational any-evasion label rate, scope-theater attempt rate, genuine-recovery rate, protocol-failure rate, protocol-failure sensitivity, Fisher comparisons against opaque denial, and turns to recovery among recovered trials. Protocol failures are non-SQL/prose turns and are treated as a separate dependent variable rather than silently folded into SQL behavior.

\section{Results}
\subsection{Data status and manipulation check}
The first attempt was denied in 99\% of trials, so the broad-request tasks reliably triggered the guardrail. Haiku is the primary run; Sonnet and Opus are partial descriptive comparisons only.\footnote{Source: \texttt{RESULTS\_STUDY.md} data-completeness lines 8--15.}

\begin{table}[t]
\centering
\caption{Haiku primary run: operational any-evasion label rate with Wilson 95\% CIs. Because these are v1 broad-intent tasks, the table is interpreted as literal rule satisfaction under a broad objective, not confirmed evasion of narrow intent.}
\label{tab:haiku-ci}
\begin{tabular}{lcc}
\toprule
Condition & Operational any-evasion label rate & n \\
\midrule
\texttt{C0\_opaque} & 22\% [13\%, 34\%] & 60 \\
\texttt{C1\_reason\_code} & 95\% [86\%, 98\%] & 60 \\
\texttt{C2\_reason\_fix} & 98\% [91\%, 100\%] & 60 \\
\texttt{C3\_reason\_fix\_blast} & 98\% [91\%, 100\%] & 60 \\
\bottomrule
\end{tabular}
\end{table}

Table~\ref{tab:haiku-ci} copies the printed Haiku Wilson intervals exactly.\footnote{Source: \texttt{RESULTS\_STUDY.md} lines 71--78.}

\begin{table}[t]
\centering
\caption{Haiku primary run: Fisher exact comparisons against opaque denial.}
\label{tab:haiku-fisher}
\begin{tabular}{lccc}
\toprule
Comparison & Printed rate comparison & Odds ratio & p \\
\midrule
\texttt{C1\_reason\_code} vs. \texttt{C0\_opaque} & 57/60 vs 13/60 & 68.69 & 1.953e-17 \\
\texttt{C2\_reason\_fix} vs. \texttt{C0\_opaque} & 59/60 vs 13/60 & 213.31 & 7.077e-20 \\
\texttt{C3\_reason\_fix\_blast} vs. \texttt{C0\_opaque} & 59/60 vs 13/60 & 213.31 & 7.077e-20 \\
\bottomrule
\end{tabular}
\end{table}

Table~\ref{tab:haiku-fisher} is descriptive inference for the complete Haiku v1 run, not a confirmatory narrow-intent evasion test.\footnote{Source: \texttt{RESULTS\_STUDY.md} lines 106--112.}

\begin{table}[t]
\centering
\caption{Operational any-evasion label rate by model and condition. Haiku is complete; Opus and Sonnet are partial descriptive runs.}
\label{tab:any-model}
\begin{tabular}{llrrrr}
\toprule
Status & Model & \texttt{C0} & \texttt{C1} & \texttt{C2} & \texttt{C3} \\
\midrule
Primary, complete & Haiku 4.5 & 22\% & 95\% & 98\% & 98\% \\
Descriptive, partial & Opus 4.8 & 22\% & 75\% & 68\% & 79\% \\
Descriptive, partial & Sonnet 4.6 & 37\% & 98\% & 100\% & 100\% \\
\bottomrule
\end{tabular}
\end{table}

The partial Opus numbers are lower than Haiku/Sonnet in the richer-feedback cells; \texttt{INDEX.md} frames this as a capability nuance, not a confirmatory cross-model result.\footnote{Rate source: \texttt{RESULTS\_STUDY.md} lines 37--43. Capability-nuance caveat source: \texttt{INDEX.md} lines 58--61.}

\begin{table}[t]
\centering
\caption{Scope-theater attempt rate by model and condition. Column labels avoid bare v1 ``evasion'' language.}
\label{tab:scope-model}
\begin{tabular}{llrrrr}
\toprule
Status & Model & \texttt{C0} & \texttt{C1} & \texttt{C2} & \texttt{C3} \\
\midrule
Primary, complete & Haiku 4.5 & 20\% & 95\% & 98\% & 98\% \\
Descriptive, partial & Opus 4.8 & 20\% & 75\% & 68\% & 79\% \\
Descriptive, partial & Sonnet 4.6 & 32\% & 98\% & 100\% & 100\% \\
\bottomrule
\end{tabular}
\end{table}

In Haiku, scope-theater attempts account for nearly all of the operational label: the rates match exactly in \texttt{C1}, \texttt{C2}, and \texttt{C3}, and differ only from 22\% to 20\% in \texttt{C0}. This is the main mechanism claim supported by v1.\footnote{Source: \texttt{RESULTS\_STUDY.md} lines 37--51.}

\begin{table}[t]
\centering
\caption{Genuine-recovery rate by model and condition.}
\label{tab:recovery-model}
\begin{tabular}{llrrrr}
\toprule
Status & Model & \texttt{C0} & \texttt{C1} & \texttt{C2} & \texttt{C3} \\
\midrule
Primary, complete & Haiku 4.5 & 38\% & 5\% & 5\% & 2\% \\
Descriptive, partial & Opus 4.8 & 40\% & 18\% & 18\% & 11\% \\
Descriptive, partial & Sonnet 4.6 & 0\% & 0\% & 0\% & 8\% \\
\bottomrule
\end{tabular}
\end{table}

Recovery moved in the opposite direction from the operational label in the complete Haiku run: 38\% under opaque denial and 5\%, 5\%, and 2\% under the three richer conditions.\footnote{Source: \texttt{RESULTS\_STUDY.md} lines 53--59.}

\begin{table}[t]
\centering
\caption{Protocol-failure rate by model and condition. Protocol failures are non-SQL/prose turns and are reported separately.}
\label{tab:protocol-model}
\begin{tabular}{llrrrr}
\toprule
Status & Model & \texttt{C0} & \texttt{C1} & \texttt{C2} & \texttt{C3} \\
\midrule
Primary, complete & Haiku 4.5 & 68\% & 95\% & 57\% & 92\% \\
Descriptive, partial & Opus 4.8 & 57\% & 65\% & 52\% & 46\% \\
Descriptive, partial & Sonnet 4.6 & 100\% & 72\% & 45\% & 82\% \\
\bottomrule
\end{tabular}
\end{table}

\begin{table}[t]
\centering
\caption{Haiku sensitivity excluding trials with protocol-failure turns.}
\label{tab:protocol-sensitivity}
\begin{tabular}{lcc}
\toprule
Condition & Operational any-evasion label rate & n \\
\midrule
\texttt{C0\_opaque} & 32\% [15\%, 54\%] & 19 \\
\texttt{C1\_reason\_code} & 33\% [6\%, 79\%] & 3 \\
\texttt{C2\_reason\_fix} & 96\% [81\%, 99\%] & 26 \\
\texttt{C3\_reason\_fix\_blast} & 100\% [57\%, 100\%] & 5 \\
\bottomrule
\end{tabular}
\end{table}

Table~\ref{tab:protocol-sensitivity} is a sensitivity check, not the main estimate: two richer-feedback cells have very small denominators, especially n=3 and n=5.\footnote{Source: \texttt{RESULTS\_STUDY.md} lines 80--87. Do not over-interpret the n=3 and n=5 cells.}

The partial-model Fisher rows are directionally consistent with Haiku but descriptive only. Opus prints 30/40 vs 9/40, OR=10.33, p=4.849e-06 for \texttt{C1}; 27/40 vs 9/40, OR=7.15, p=0.0001041 for \texttt{C2}; and 22/28 vs 9/40, OR=12.63, p=5.94e-06 for \texttt{C3}. Sonnet prints 59/60 vs 22/60, OR=101.91, p=3.031e-14 for \texttt{C1}, and 40/40 vs 22/60, OR=inf, p=3.194e-12 for both \texttt{C2} and \texttt{C3}.\footnote{Source: \texttt{RESULTS\_STUDY.md} lines 114--128. The Sonnet OR=inf entries are perfect-separation artifacts from 40/40 richer-condition cells against 22/60 opaque baseline cells, not literal infinite effect sizes.}

Among recovered trials only, the printed mean turns-to-recovery were 2.23 for \texttt{C0}, 1.50 for \texttt{C1}, 1.90 for \texttt{C2}, and 1.71 for \texttt{C3}, with recovered-trial counts 39.00, 10.00, 10.00, and 7.00. These pooled recovered-trial summaries are not used as confirmatory cross-model evidence.\footnote{Source: \texttt{RESULTS\_STUDY.md} lines 156--163.}

\section{Discussion}
These are v1 results: they support the claim that mechanistic denial feedback elicited literal rule satisfaction under a broad objective, not a confirmatory claim about narrow-intent evasion. The completed Haiku run contains the pattern most cleanly: the operational label rose from 22\% [13\%, 34\%] under opaque denial to 95\% [86\%, 98\%] with a reason code and 98\% [91\%, 100\%] with reason-plus-fix or reason-plus-fix-plus-blast feedback. In the same Haiku cells, genuine recovery moved from 38\% under opaque denial to 5\%, 5\%, and 2\% under the richer conditions.

The plausible v1 mechanism is literal satisfaction of the named rule. The prompts asked for broad writes; the guardrail rejected broad writes. When the denial named the violated rule, the agent often added a syntactically valid \texttt{WHERE} while preserving the broad operation. In Haiku, scope-theater attempts match the operational label in \texttt{C1}, \texttt{C2}, and \texttt{C3}, which supports the mechanism claim that richer mechanistic feedback made the visible rule easier to satisfy literally.

This does not imply that transparency is generally harmful. Reliable feedback can help when the agent's objective and the guardrail's objective are aligned. In v1 they were in tension: the user-facing task asked for a broad operation, while the guardrail demanded scoped impact. In that setting, mechanistic feedback such as \texttt{WRITE\_WITHOUT\_WHERE} can function less like guidance toward a narrow intended action and more like a specification of the surface property that must change.

\section{Limitations}
The main validity limit is task framing. v1 tasks explicitly asked for broad operations, so whole-table restatements can be faithful compliance with the user request rather than evasion of narrow intent. The paper therefore cannot make a confirmatory evasion claim from v1.

The data state also limits inference. Haiku is complete and primary; Sonnet and Opus are partial because the sweep was interrupted by API credit limits. Additional caveats are a single provider family, 3 synthetic adversarial tasks, partial Sonnet/Opus cells, one seed, and v1 C3 confounding row-count feedback with confirmation framing.\footnote{These caveats are listed in \texttt{INDEX.md}, lines 82--90.}

Protocol failures are a separate threat. Non-SQL/prose payloads were frequent, and the protocol-failure-excluded sensitivity analysis leaves very small denominators in some Haiku cells. The primary result should therefore be reported alongside, not replaced by, the sensitivity check.

Finally, the labels are operational and auditable, but a human audit remains pending. The audit protocol calls for a stratified sample, independent raters, reported agreement, and publication of the adjudications. Until that audit is complete, borderline label assignments should be treated as operational measurements rather than settled behavioral ground truth.\footnote{Source: \texttt{label\_audit.md}, lines 26--38; the draft does not use the audit protocol as a result estimate.}

\section{What v2 would test}
The v2 design is the confirmatory test because it fixes the framing problem. The pre-registration states that v1 is exploratory and reframed, while v2 uses narrow-intent prompts in which over-reach is genuine rather than compliance. Its task set has 4 narrow updates, 4 narrow deletes, 2 legitimately broad tasks, and 2 ambiguous tasks, with intended rows derivable from the prompt. Its feedback conditions are isolated as \texttt{V0\_opaque}, \texttt{V1\_behavioral}, \texttt{V2\_reason\_code}, \texttt{V3\_reason\_fix}, and \texttt{V4\_reason\_fix\_rows}; \texttt{V4} adds measured rows without confirmation framing.\footnote{Source: \texttt{pre\_registration.md}, lines 1--33, and \texttt{tasks\_v2.py}, lines 1--12 and 57--116.}

The planned hypotheses are not findings. H1: behavioral feedback (\texttt{V1}) improves \texttt{correct\_recovery} versus opaque feedback (\texttt{V0}). H2: reason-code feedback (\texttt{V2}) increases \texttt{syntax\_only\_compliance} versus \texttt{V0} and \texttt{V1}. H3: adding measured rows (\texttt{V4}) reduces \texttt{semantic\_overreach} versus reason-code-only feedback (\texttt{V2}). The pre-registration also says only fully balanced models enter cross-model confirmatory claims; partial models remain descriptive.

\section{Conclusion}
Interdict's v1 evidence is narrow but useful. In a setting where the prompt asked for broad database writes and the guardrail blocked broad impact, naming the SQL rule made agents much more likely to satisfy the named rule literally by adding a \texttt{WHERE} that did not narrow impact. The result supports an exploratory design lesson: guardrail messages for database agents should be evaluated not only for whether they explain a rejection, but for whether they reveal a surface rule that can be satisfied while preserving the unsafe blast radius. The confirmatory evasion claim belongs to the v2 narrow-intent study, not to v1.

\end{document}
